\documentclass[conference]{IEEEtran}
\IEEEoverridecommandlockouts
% The preceding line is only needed to identify funding in the first footnote. If that is unneeded, please comment it out.
\usepackage{cite}
\usepackage{amsmath,amssymb,amsfonts}
\usepackage{algorithmic}
\usepackage{graphicx}
\usepackage{textcomp}
\usepackage{xcolor}
\usepackage{caption}
\usepackage{subcaption}
\usepackage{url}
\def\BibTeX{{\rm B\kern-.05em{\sc i\kern-.025em b}\kern-.08em
    T\kern-.1667em\lower.7ex\hbox{E}\kern-.125emX}}
\begin{document}

\title{Masters Project Final Report}



\author{\IEEEauthorblockN{Khoa Le}
\IEEEauthorblockA{\textit{University of Georgia}}
\textit{Hetergeneous Robotics Research Lab}\\
Athens, Georgia \\
kl58277@uga.edu}

\maketitle

\begin{abstract}
This project developed and deployed a Rover Robotics platform for autonomous poultry-farm operations, with the goal of serving multiple research projects on chicken farm monitoring. For example, measuring chicken's travel distance based on the movement of the robot, or scooping up dead chickens. Modern poultry houses present a dense, dynamic, and sensitive environment where manual data collection or experiments are labor-intensive and inefficient. To address this , I worked on a mobile Rover Robot with LiDar-driven perception, ROS2 navigation, and a variety of configurable modules to support a range of operational tasks. Beyond locomotion and sensing, my central research objective was replicating and extending the Collaborative Direction of Arrival (CDOA)–based RSSI localization method proposed by Latif $\And$ Parasuraman \cite{b1}, evaluating its practicality for real-world field robotics. The rover employs LiDAR perception, ROS2 navigation, and a configurable sensing stack while incorporating an RSSI-only localization module designed to operate in GPS-denied agricultural environments. The CDOA algorithm leverages cooperative RSSI measurements from distributed wireless nodes to estimate the robot’s direction of arrival and refine its position using Bayesian inference (EM and PF). This final report presents the design, implementation, and evaluation of a wireless indoor localization system for poultry-farm robotics using ROS 2 Jazzy. 
To date, testing has been conducted in the lab environments and an indoor testbed designed to approximate key features of a poultry house. My future work will involve full deployment and validation inside an operational poultry farm to assess robustness, animal interaction dynamics, and real-world research utility. This ongoing effort aims to establish an extensible robotic framework for precision indoor poultry navigation.
\end{abstract}

\begin{IEEEkeywords}
Rover Robotics, localization, navigation, poultry
\end{IEEEkeywords}

\section{Introduction}
The poultry industry relies on continuous monitoring of bird health, behavior, and environmental conditions to maintain welfare standards and operational efficiency. However, modern commercial poultry houses present substantial challenges for manual observation. Large flock sizes, uneven terrains, and large scale environment make traditional data collection methods labor-intensive and is extremely difficult. As a result, researchers continue to seek automated methods, such as integrating robots that can support measurement of flock activity, to reduce the amount of work needed to be done manually. 

Indoor localization is a foundational capability for mobile robots operating in GPS-denied environments such as commercial poultry houses. These environments are  cluttered and noisy where LiDAR and camera-based localization can degrade, and where full SLAM pipelines may fail due to occlusions or rapidly changing scene geometry. For agricultural researchers, these limitations constrain the use of mobile robots for tasks such as monitoring flock activity, estimating chicken travel distance, and conducting automated welfare inspections.

\begin{figure}[htbp]
    \centering
    \includegraphics[width=0.48\textwidth]{poultry.jpeg} 
    \caption{A full-scale chicken farm in Georgia where the experiment will be conducted}
    \label{fig:poultry}
\end{figure}

Recent work in wireless-signal–based localization offers a promising alternative. Latif and Parasuraman’s CDOA framework demonstrated that Received Signal Strength Indicator (RSSI) measurements from a small number of static wireless sensor nodes can be used to estimate the Direction of Arrival (DOA) of a mobile node, and that combining CDOA with Bayesian methods such as Expectation Maximization (EM) or Particle Filtering (PF) enables accurate indoor localization without fingerprinting or expensive sensors \cite{b1}. This approach is particularly efficient for agricultural robotics, where simplicity, low cost, and robustness to visual noise are critical.

My project aims to adapt, replicate, and scale the CDOA methodology for a Rover Robotics platform intended for poultry farm operations. The rover is designed as a multi-purpose research tool capable of estimating flock mobility by using its own movement as a surrogate for bird travel distance, performing mortality retrieval tasks through mechanical attachments, and supporting modular environmental sensing. To enable reliable navigation and data collection, I am implementing a wireless localization system based on collaborative RSSI sampling across distributed access points. My goal is to assess whether the CDOA method can perform effectively in the challenging conditions of a real poultry house.

To date, experiments have been conducted in the lab settings and in an indoor testbed. These evaluations have focused on verifying the reproducibility of the CDOA algorithm and integrating the localization output with the rover’s ROS2 navigation stack. Future work includes deploying the system inside an active poultry farm to examine localization stability under dynamic flock movement. Ultimately, this project aims to bridge wireless signal–based localization research with real-world agricultural robotics, establishing a pathway for low-cost, scalable robot navigation in commercial poultry production.

\section{Literature Review}

Indoor localization in GPS-denied environments remains a core challenge in mobile robotics, particularly in settings where visual sensors and LiDAR can fail due to rapid change in the environment geometry due to chicken's movements. Latif $\And$ Parasuraman (2023) \cite{b1} present one of the most promising recent advances in wireless-signal–based localization: the Collaborative Direction of Arrival (CDOA) method, which enables instantaneous localization of mobile robotic nodes using only RSSI data from a small number of static wireless sensors (4 count).

The key contribution of their work is calculating the RSSI gradients, which are measured collaboratively across multiple wireless sensor nodes. This can be used to compute the estimate of the Direction of Arrival without requiring fingerprinting. This is achieved through the central finite-difference method, which uses spatially distributed RSSI readings to approximate the signal gradient along the x and y axes. The resulting gradient is then converted into a CDOA angle, providing a bearing estimate for the mobile robot.

A relevant aspect of this work is its focus on resource-constrained environment and low-cost wireless infrastructure. Unlike vision-based or LiDAR-based localization methods, CDOA does not rely on lighting or GPS. This method requires only simple wireless nodes and RSSI measurements, making it very efficient and robust in difficult environments. 

Overall, the CDOA work by Latif and Parasuraman forms a strong foundation for robotics research seeking low-cost, scalable indoor localization without reliance on vision or LiDAR. Their findings suggest that wireless-only localization can outperform classical RSSI techniques and provide practical accuracy for mobile robot navigation. This is perfect for applications such as poultry-farm robotics, where the dynamic factors degrade visual and GPS-based methods.

\section{Design and Implementation}
The primary objective of this project is to replicated and extend the Collaborative Direction of Arrival (CDOA) indoor localization method proposed by Latif $\And$ Parasuraman \cite{b1} and integate it into a Rover Robot. The following subsections describe the system architecture, hardware setup, and software implementation. A total of nine wireless nodes are deployed to enable large and accurate coverage for RSSI measurement. 

\subsection{System Architecture Overview}
The final system contains three components:
\begin{itemize}
\item A mobile robot platform running ROS 2 Jazzy with:
\begin{itemize}
\item Robot odometry input from wheel odometry topic
\item Velocity input from command velocity topic
\item PF localization node for wireless-based pose estimation
\item Teleoperation and launch-based runtime control
\end{itemize}
\item A wireless sensing network composed of four static Raspberry Pi 5 nodes, each publishing wireless link measurements.
\item A network-analysis and localization software stack where:
\begin{itemize}
\item Wireless quality messages are published from each static node
\item Network diagnostics (errors, utilization, and delay) can run in parallel
\item The localization node subscribes to four wireless streams and robot motion streams
\item The PF loop estimates robot position from RSSI-derived directional information and motion update
\end{itemize}
\end{itemize}

The implementation uses the updated ROS 2 message interfaces for wireless and network metrics. Wireless measurements include RSSI, link quality indicator (LQI), noise, interface metadata, and link status. This enables the localization algorithm to run while also monitoring communication quality during experiments.

\subsection{Visualization}
\begin{figure}[htbp]
    \centering
    \includegraphics[width=0.48\textwidth]{setup1.png} 
    \caption{Setup 1}
    \label{fig:setup1}
\end{figure}
\begin{figure}[htbp]
    \centering
    \includegraphics[width=0.48\textwidth]{setup2.png} 
    \caption{Setup 2}
    \label{fig:setup2}
\end{figure}

The experimental setup consists of a teleoperated mobile robot positioned at the center of a distributed wireless sensor network composed of eight static nodes arranged around the workspace. Each static node measures the Wi-Fi RSSI broadcast from the access point to enable signal sampling for CDOA-based gradient estimation. The two configurations illustrated in Figures 1 and 2 represent variations in node spacing and geometric distribution designed to evaluate the sensitivity of the CDOA method to environmental layout.  These controlled configurations provide a structured system for analyzing the impact of node placement on overall system performance within a ROS2 Jazzy framework.

\subsection{Code Implementation}
To support the CDOA localization algorithm, modifications were required to the existing \texttt{ros2-network-analysis} GitHub repository. The original package provides basic network diagnostics such as RSSI, link quality, packet errors, and throughput; however, a lot of the code were not compatible with the new ROS2 Jazzy package. Therefore, several components of the repository were modified, specifically the \texttt{network\_errors.py} file, which included deprecated syntax. 

\begin{figure}[htbp]
    \centering
    \includegraphics[width=0.48\textwidth]{error3.png} 
    \caption{Original \texttt{network\_errors.py} File }
    \label{fig:error3}
\end{figure}
Originally, this file contains an errors on line such as \texttt{msg.retransmits}. The problem was that this line is read from standard output and casted on an integer. However, the data type of read() is \textbf{byte}. Therfore, I have made changes throughout the codes to first decode it then cast after. Figure 4. shows the modified code. 
\begin{figure}[htbp]
    \centering
    \includegraphics[width=0.48\textwidth]{error2.png} 
    \caption{Modified \texttt{network\_errors.py} File }
    \label{fig:error2}
\end{figure}

Another problem with the localization was that the area size and the node boundary were set incorrectly. In the Figure 5., 
the area size and node positions were measured in meters. I have realized that I need to modify these to match the real-world measurements; this original code was used on a test-bed.
\begin{figure}[htbp]
    \centering
    \includegraphics[width=0.48\textwidth]{error4.png} 
    \caption{Original \texttt{PFLocalization\_tabletop.py} File }
    \label{fig:error4}
\end{figure}

For this final project, the localization and network measurement modules were implemented and migrated into ROS 2 Jazzy. The localization node subscribes to four wireless quality topics, wheel odometry, and command velocity. In each update cycle, the node stores incoming measurements, computes directional information from RSSI gradients, and performs PF-based state estimation.

To support localization and runtime analysis, the following software components were integrated:
\begin{itemize}
\item Wireless-quality publisher for RSSI, LQI, noise, and link status
\item Network-error publisher for retransmissions, segment errors, and interface error counters
\item Link-utilization publisher for transmit/receive throughput and packet-level rates
\item Ping service-client pair for network-delay estimation and testing purposes
\end{itemize}

\subsection{Localization Runtime Pipeline}
The localization implementation follows a measurement-update and motion-update sequence:
\begin{enumerate}
\item Collect synchronized robot and wireless observations from ROS 2 topics.
\item Form a four-node RSSI vector and compute directional gradient terms.
\item Generate expected RSSI values for particle hypotheses using a wireless propagation model.
\item Compute likelihood weights from RSSI-direction error statistics.
\item Select the best weighted hypothesis and combine it with robot trajectory context.
\item Apply a motion-model update using linear and angular velocity input.
\item Resample particles in a local region around the current estimate.
\end{enumerate}
\begin{figure}[htbp]
    \centering
    \includegraphics[width=0.48\textwidth]{figure3.png} 
    \caption{The system architecture of the CDOA-based localization algorithm, which uses Particle Filter and Estimation Maximization. This figure is from paper \cite{b1}}
    \label{fig:fig3}
\end{figure}

The implementation also plots estimated and reference trajectories online, allowing immediate visual verification of localization behavior during experiments (see Fig. 7).

\begin{figure}[htbp]
    \centering
    \includegraphics[width=0.48\textwidth]{error.jpg} 
    \caption{My Localization Output}
    \label{fig:localization_ouput}
\end{figure}

\subsection{Implementation Notes and Constraints}
The current implementation is functional for indoor testing, but several constraints were observed:
\begin{itemize}
\item Topic naming must be aligned across distributed nodes and robot runtime.
\item Localization behavior is sensitive to node geometry and signal variability.
\item Several localization parameters are currently fixed in code for testbed operation.
\item You must set up the boundary of the Raspberry Pi nodes correctly. If not, the algorithm will be very off. 
\end{itemize}

Despite these constraints, the completed ROS 2 pipeline provides a reproducible baseline for wireless localization experiments and supports direct extension toward larger poultry-house deployment

\section{Discussion}
By the end of this project, the Rover Robotics platform had been configured for autonomous navigation in ROS 2, and the wireless localization stack was integrated into the robot workflow. The full runtime procedure for building, launching, and running the system is documented in the repository, so future users can follow the GitHub instructions to reproduce the complete setup from code to execution. This makes the repository a practical deployment guide as well as a research archive for poultry-robot experiments.

The main limitation of the Rover robots is reliability. In most runs the system works as intended, but occasionally the robot crashes or becomes unstable, so a full system reset is sometimes required before continuing experiments. This means that the user has to click the power button to turn off. This makes the platform usable for testing and development, but it also shows that longer autonomous sessions still require hardware and software stabilization.

The Husarion Rosbot introduces a different limitation because it runs ROS 2 Foxy instead of ROS 2 Jazzy. To inspect its runtime visualization tools such as RViz2, I recommend users to SSH into the robot, but this adds setup time and slows down debugging. In addition, subscribing to its topics requires using the same ROS distribution, so cross-distribution integration is not straightforward. These constraints are important for future researchers who want to reuse the robot in a compatible ROS environment.
\begin{figure}[htbp]
    \centering
    \includegraphics[width=0.48\textwidth]{rosbot.jpg} 
    \caption{Husarion Rosbot 2 Pro}
    \label{fig:rosbot}
\end{figure}
\begin{figure}[htbp]
    \centering
    \includegraphics[width=0.48\textwidth]{rover.jpg} 
    \caption{Rover Robot Max}
    \label{fig:rover}
\end{figure}
Overall, this repository was created to serve as a comprehensive reference for future researchers who want to use these robots for poultry operations, indoor mobility experiments, and wireless localization studies. It documents the ROS 2 workflow, launch procedure, and system integration steps needed to support repeatable experimentation in agricultural robotics.
Repository can be found at: \textbf{\url{https://github.com/lenhatdangkhoa/cdoa-ros2}}
\section{Future Work}
Future work will focus on improving localization robustness and expanding the system to larger-scale poultry-house deployments. The first priority is to increase reliability on the Rover platform by reducing crash rate and improving automated recovery behavior so that experiments can run for longer periods without manual intervention. In parallel, the wireless localization pipeline should be evaluated with more trajectories, more trials, and more controlled node geometries to obtain stronger quantitative performance results.

A second direction is to improve integration between localization, odometry, and IMU data so that the robot can maintain more stable pose estimates under noisy wireless conditions. The current implementation already provides a functional ROS 2 baseline, but higher numbers of nodes would make the system more suitable for autonomous navigation and longer-duration large field experiments. For the Husarion Rosbot, future work should also focus on simplifying cross-distribution access so that visualization, topic subscription, and debugging can be performed more efficiently.

Ultimately, the goal is to turn this repository into a complete practical guide for poultry-robotics research, with reusable launch procedures, clearer deployment instructions, and a more robust localization and navigation stack for future researchers.

\section*{Acknowledgment}
I would like to express my sincere gratitude to my Professor Ramviyas Nattanmai Parasuraman for his guidance and feedback throughout the development of this whole project. I would also like to thank my friend Daniel Laij for his support and collaboration during the debugging phases. Their contributions have been invaluable to my progress for this project. 
\section{The use of Generative AI tool}
\textbf{I have used ChatGPT to help me polish some of my writing to make it sounds for professional. I did not use it to generate this report. }

\begin{thebibliography}{00}
\bibitem{b1} E. Latif and R. Parasuraman, ``Instantaneous Wireless Robotic Node Localization Using Collaborative Direction of Arrival,'' *IEEE Internet of Things Journal*, accepted for publication, 2023. Available: https://arxiv.org/abs/2307.01956.


\end{thebibliography}

\end{document}
